% This file is intended to be included from the main manuscript with \input.
% It uses the notation and theorem/algorithm environments defined in the
% manuscript preamble.  The first section can replace the current subsection
% labeled sec:unbiased_estimation.  The second section may be placed in the
% methodological appendix.

\section{Population-level metric estimation under adaptive censoring}
\label{sec:unbiased_estimation}

Beyond constructing an LPB for an individual prompt, an evaluator may wish to
summarize the behavior of the target model over an entire benchmark.  For
example, in a safety application, we may ask what fraction of prompts elicit
an unsafe response within a horizon of $\maxl$ interactions.  A naive answer
based only on the observed unsafe events is generally biased: whenever the
budget allocator stops a conversation early, an unsafe event that would have
occurred later is recorded as if it never occurred.  This bias is not a
minor technicality.  A policy that saves the most compute by stopping
apparently safe conversations early can also create the largest distortion if
the resulting censoring is ignored.

In this section, we show that the acquisition propensities recorded by our
allocation policy remove this bias exactly.  We first present the result for a
general event-supported metric, then specialize it to the unsafe event rate.
The same argument applies to successful-event rates in utility evaluation.

\subsection{Resolved outcomes and pathwise acquisition propensities}

Fix the maximal conversational horizon $\maxl$.  For each $i\in\calI$, let
\begin{equation}
    A_i := \ind{T_i\leq \maxl}
\end{equation}
denote the event that the conversation produces an unsafe response by the
horizon.  At interaction $t$, the budget allocator observes the currently
available history $H_i(t)$ and chooses a continuation probability
$P_i(t)\in[0,1]$.  It then draws
\begin{equation}
    Z_i(t)\sim\operatorname{Ber}(P_i(t)).
\end{equation}
The policy is \emph{causal}: $P_i(t)$ may depend on $X_i$, $H_i(t)$, the
training data $\Dtrain$, and an independent policy-fitting split, but not on a
future response or on the unobserved value of $T_i$.

For an unsafe event occurring at $T_i\leq\maxl$, define
\begin{equation}
    R_i := \prod_{t=1}^{T_i} Z_i(t),
    \qquad
    \pi_i := \prod_{t=1}^{T_i} P_i(t).
    \label{eq:metric_event_inclusion}
\end{equation}
Thus, $R_i=1$ precisely when the unsafe event is observed, and $\pi_i$ is its
pathwise inclusion probability.  Although $P_i(t)$ is adaptive, the product
in~\eqref{eq:metric_event_inclusion} is known: each factor is selected by the
auditor before the corresponding Bernoulli draw.  We use the harmless
conventions $R_i=0$ and $A_iR_i/\pi_i=0$ when $A_i=0$.

More generally, let $g:\{1,\ldots,\maxl\}\rightarrow\R$ and define the
event-supported individual contribution
\begin{equation}
    G_i := g(T_i)A_i.
\end{equation}
Examples include $G_i=A_i$ for the unsafe event rate and
$G_i=T_iA_i$ for the event-time numerator introduced below.  The finite-
benchmark estimand and its Horvitz--Thompson estimator are
\begin{equation}
    \mu_G := \frac{1}{N}\sum_{i\in\calI}G_i,
    \qquad
    \widehat\mu_G := \frac{1}{N}\sum_{i\in\calI}
    \frac{R_iG_i}{\pi_i}.
    \label{eq:metric_ht_estimator}
\end{equation}
Fully observed policy-fitting or budget-control points are included in this
notation by setting $R_i=\pi_i=1$.

\begin{assumption}[Causal acquisition and positivity]
\label{assump:metric_causal_positive}
Conditional on the complete latent trajectories, $\Dtrain$, and any data used
to fit the acquisition policy, the Bernoulli acquisition variables are
independent across prompts and are drawn according to their recorded
conditional probabilities.  Moreover, there exists $\eps>0$ such that every
target event has $\pi_i\geq\eps$ almost surely.
\end{assumption}

Assumption~\ref{assump:metric_causal_positive} does \emph{not} require a
conditionally independent censoring model of the form $C_i\indep T_i\mid
X_i$.  The continuation probability may react to the entire observed
conversation history.  What is required is that future information is not
used before it is observed and that the probability generated by the policy
is recorded exactly.

\begin{theorem}[Design validity of adaptive metric estimation]
\label{thm:metric_ht_validity}
Suppose Assumption~\ref{assump:metric_causal_positive} holds.  Conditional on
the complete benchmark trajectories and on the fitted acquisition policy,
the estimator in~\eqref{eq:metric_ht_estimator} is unbiased:
\begin{equation}
    \E\!\left[\widehat\mu_G\mid
    \{X_i,H_i,T_i\}_{i\in\calI},\Dtrain,\mathcal{M}\right]
    = \mu_G,
    \label{eq:metric_ht_conditional_unbiasedness}
\end{equation}
where $\mathcal{M}$ denotes all randomness and data used to fit the frozen
acquisition policy.  Consequently, it is also unbiased after averaging over
the random benchmark and policy-fitting split.
\end{theorem}

\begin{proof}
Fix an index with $A_i=1$.  Conditional on its complete trajectory and the
fitted policy, the sequential tower property gives
\begin{equation}
    \E[R_i\mid X_i,H_i,T_i,\Dtrain,\mathcal{M}]
    = \prod_{t=1}^{T_i}P_i(t)=\pi_i.
\end{equation}
Hence
\begin{equation}
    \E\!\left[\frac{R_iG_i}{\pi_i}\middle|
    X_i,H_i,T_i,\Dtrain,\mathcal{M}\right]=G_i.
\end{equation}
When $A_i=0$, both sides are zero.  Summing over $i$ proves
\eqref{eq:metric_ht_conditional_unbiasedness}; taking another expectation
proves the marginal statement.
\end{proof}

For the unsafe event rate, $G_i=A_i$, and Theorem~\ref{thm:metric_ht_validity}
yields the estimator
\begin{equation}
    \widehat{\ttuer}
    :=\frac{1}{N}\sum_{i\in\calI}\frac{R_iA_i}{\pi_i},
    \qquad
    \ttuer:=\frac{1}{N}\sum_{i\in\calI}A_i.
    \label{eq:uer_ht}
\end{equation}
Notice that a safe conversation need not be run to $\maxl$ in order to
contribute a zero to~\eqref{eq:uer_ht}.  This is the main computational
advantage of estimating an event-supported metric: compute is needed to
reveal unsafe events, not to certify every safe trajectory individually.
Operationally, the evaluator sums the weighted contributions of the observed
unsafe events; every other row contributes zero, whether it is truly safe or
was censored before a later event.

\subsection{Exact conditional variance and the $A$-weighted objective}

Unbiasedness alone is not enough.  If an unsafe event is observed with an
extremely small probability, its contribution is multiplied by a very large
inverse propensity, producing an estimator that is valid but unstable.  The
following result identifies the allocation objective that controls this
instability.

\begin{theorem}[Conditional acquisition variance]
\label{thm:metric_ht_variance}
Under the conditions of Theorem~\ref{thm:metric_ht_validity}, and conditional
on the complete benchmark trajectories and fitted policy,
\begin{equation}
    \operatorname{Var}(\widehat\mu_G)
    =\frac{1}{N^2}\sum_{i\in\calI}
    G_i^2\left(\frac{1}{\pi_i}-1\right).
    \label{eq:general_metric_variance}
\end{equation}
In particular,
\begin{equation}
    \operatorname{Var}(\widehat{\ttuer})
    =\frac{1}{N^2}\sum_{i\in\calI}
    A_i\left(\frac{1}{\pi_i}-1\right).
    \label{eq:uer_variance}
\end{equation}
\end{theorem}

\begin{proof}
Conditional on the complete trajectory and policy, $R_i$ is Bernoulli with
mean $\pi_i$.  Therefore,
\begin{align}
    \E\!\left[\left(\frac{R_iG_i}{\pi_i}\right)^2\right]
    &=\frac{G_i^2}{\pi_i^2}\E[R_i]
      =\frac{G_i^2}{\pi_i},\\
    \operatorname{Var}\!\left(\frac{R_iG_i}{\pi_i}\right)
    &=G_i^2\left(\frac{1}{\pi_i}-1\right).
\end{align}
The acquisition draws are independent across prompts conditional on the
fitted policy, so the variances add.  Finally, $A_i^2=A_i$ gives
\eqref{eq:uer_variance}.
\end{proof}

Theorem~\ref{thm:metric_ht_variance} explains why minimizing the ordinary
mean weight $N^{-1}\sum_i1/\pi_i$ can be inefficient for metric estimation.
That objective spends budget to increase the inclusion probability of every
prompt, including prompts with $A_i=0$, whose contribution to the variance in
\eqref{eq:uer_variance} is exactly zero.  The variance-aligned objective is
instead
\begin{equation}
    \frac{1}{N}\sum_{i\in\calI}
    A_i\left(\frac{1}{\pi_i}-1\right),
    \label{eq:a_weighted_metric_objective}
\end{equation}
which we call the \emph{$A$-weighted objective}.  The additive term
$-N^{-1}\sum_iA_i$ does not depend on the allocation and may be omitted by
the optimizer.

The oracle $A$-weighted objective cannot be optimized directly because the
unobserved $A_i$'s are precisely what the evaluation seeks to discover.  The
allocation methods below differ in how they estimate this objective without
using future information: hard DAPRO estimates it from fully observed Phase-I
endpoints, Generalized DAPRO replaces endpoints by causal prefix hazards, and
Metric-optimal PMF integrates over the predictive distribution available at
the initial prompt.

An unbiased estimator of the conditional variance in~\eqref{eq:uer_variance},
which uses only actually observed unsafe events, is
\begin{equation}
    \widehat V_{\mathrm{acq}}
    :=\frac{1}{N^2}\sum_{i\in\calI}
    R_iA_i\frac{1-\pi_i}{\pi_i^2}.
    \label{eq:observed_variance_estimator}
\end{equation}
Indeed, the conditional expectation of the $i$-th summand is
$A_i(1/\pi_i-1)$.  This diagnostic is useful for validating an
implementation, whereas the variance reported in our experiments is the
empirical variance of the final metric estimates across random
calibration--test splits.

\subsection{Acquisition variance versus split variance}

Our experiments repeatedly draw a calibration subset from a fixed
calibration--test pool.  This introduces a second source of randomness that
is present even if every selected calibration trajectory is observed in
full.  Let $\mathcal{I}_S$ denote a random calibration subset of size $n$ from
a fixed pool of size $N_{\mathrm{all}}$, and let
$\mu_S=n^{-1}\sum_{i\in\mathcal{I}_S}A_i$.  The law of total variance gives
\begin{equation}
\begin{split}
    \operatorname{Var}(\widehat{\ttuer})
    ={}&\operatorname{Var}_S(\mu_S)\\
    &+\E_S\!\left[
    \operatorname{Var}_{\mathrm{acq}}
    (\widehat{\ttuer}\mid\mathcal{I}_S)
    \right].
    \label{eq:split_acquisition_variance}
\end{split}
\end{equation}
For simple random sampling without replacement,
\begin{equation}
    \operatorname{Var}_S(\mu_S)
    =\left(1-\frac{n}{N_{\mathrm{all}}}\right)
      \frac{S_A^2}{n}
    =\frac{N_{\mathrm{all}}-n}{n(N_{\mathrm{all}}-1)}
      \mu_{\mathrm{all}}(1-\mu_{\mathrm{all}}),
    \label{eq:split_variance_binary}
\end{equation}
where $S_A^2$ is the finite-population variance and
$\mu_{\mathrm{all}}=N_{\mathrm{all}}^{-1}\sum_iA_i$.

Equation~\eqref{eq:split_acquisition_variance} also clarifies the role of the
two infinite-budget references used in our experiments.  Full observation of
the calibration subset sets the second term to zero but retains the split
variance in~\eqref{eq:split_variance_binary}.  Full observation of the fixed
calibration--test union has zero variance because that union does not change
across repetitions.  We use the latter as the fixed oracle value and the
former as the irreducible split-variance reference.

\subsection{Event-time metrics and ratio estimators}

The code also reports a restricted mean time-to-unsafe, denoted by
$\ttrmttu$.  More precisely, the implemented quantity is the mean event time
among prompts that experience an unsafe event by $\maxl$:
\begin{equation}
    r_{\maxl}
    :=\E[T\mid T\leq\maxl]
    =\frac{\nu_{\maxl}}{p_{\maxl}},
    \qquad
    \nu_{\maxl}:=\E[T\ind{T\leq\maxl}],
    \quad
    p_{\maxl}:=\E[\ind{T\leq\maxl}].
    \label{eq:rmttu_ratio_target}
\end{equation}
We estimate the numerator and denominator with two HT estimators,
\begin{equation}
    \widehat\nu_{\maxl}
    =\frac{1}{N}\sum_i\frac{R_iT_iA_i}{\pi_i},
    \qquad
    \widehat p_{\maxl}
    =\frac{1}{N}\sum_i\frac{R_iA_i}{\pi_i},
\end{equation}
and report
\begin{equation}
    \widehat r_{\maxl}
    =\frac{\widehat\nu_{\maxl}}{\widehat p_{\maxl}}.
    \label{eq:rmttu_hajek}
\end{equation}
Both totals are exactly unbiased.  Their ratio is a H\'ajek-type ratio
estimator: it is consistent and asymptotically unbiased under the usual
non-degeneracy conditions, but it is not exactly unbiased at finite $N$.
This distinction is important when interpreting small-sample experiments.
The leading variance is governed by the influence contribution
\begin{equation}
    \phi_i=\frac{A_i(T_i-r_{\maxl})}{p_{\maxl}},
\end{equation}
and therefore a policy designed specifically for~\eqref{eq:rmttu_hajek}
would replace $A_i$ in~\eqref{eq:a_weighted_metric_objective} by
$\phi_i^2$.  In the present experiments, the acquisition policies are tuned
for the primary $\ttuer$ target; $\ttrmttu$ is reported as a secondary
metric.

We also note that the classical restricted mean survival time
$\E[\min(T,\maxl)]$ is a different estimand.  Estimating that quantity
requires resolving both observed events and conversations known to survive
through $\maxl$.


\section{A unified DAPRO framework for predictive bounds and metric estimation}
\label{sec:generalized_dapro}

We next present a unified view of the deployable DAPRO variants used in this
work.  The central message is that LPB construction and metric estimation do
not require unrelated allocation algorithms.  Both tasks reduce to choosing
cumulative reach probabilities under the same sequential budget constraint;
they differ in the target-specific coefficients placed in the objective and,
in the initial-PMF method, in the policy class available before the first
interaction.

\subsection{The common event-mass/cost-mass problem}

Let
\begin{equation}
    \rho_i(t):=\prod_{s=1}^{t}P_i(s)
    \label{eq:cumulative_reach_probability}
\end{equation}
be the probability that the policy acquires interaction $t$ of prompt $i$.
For a fixed target, associate each prompt--time pair with two nonnegative
coefficients:
\begin{itemize}
    \item $a_i(t)$, the event or squared-influence mass whose inverse
    propensity contributes to the target estimator's variance; and
    \item $d_i(t)$, the probability or empirical indicator that interaction
    $t$ is active and therefore consumes one unit of budget when reached.
\end{itemize}
The generic design problem is
\begin{equation}
\label{eq:generalized_dapro_problem}
\begin{aligned}
    \min_{\rho\in\mathcal{P}}\quad
    &\frac{1}{n}\sum_{i=1}^{n}\sum_{t=1}^{\maxl}
      a_i(t)\left(\frac{1}{\rho_i(t)}-1\right)\\
    \text{s.t.}\quad
    &\frac{1}{n}\sum_{i=1}^{n}\sum_{t=1}^{\maxl}
      d_i(t)\rho_i(t)\leq\bar B_{\mathrm{fit}}.
\end{aligned}
\end{equation}
Here, $\mathcal{P}$ is the chosen class of causal policies.  The constant
terms $-a_i(t)$ may be removed without changing the optimizer.  The method
called \emph{Target-$A$ DAPRO}, the stabilized \emph{Definitive DAPRO}, the
history-adaptive \emph{Generalized DAPRO}, and the pre-run
\emph{Metric-optimal PMF} method are all instances of
\eqref{eq:generalized_dapro_problem}.

This representation is useful because it separates three choices that should
not be conflated:
\begin{enumerate}
    \item the statistical target, which determines $a_i(t)$;
    \item the coefficient estimator, e.g., a hard observed endpoint or a
    model-integrated soft event mass; and
    \item the policy class $\mathcal{P}$, e.g., a history-adaptive score-bin
    policy or a schedule committed before the first interaction.
\end{enumerate}
Budget control is a fourth, orthogonal choice: the same policy shape may be
deployed using a projection assumption or selected using an independent CRC
fold.

\subsection{The policy-fitting split and causal scores}

For the history-adaptive variants, randomly split $\calI$ into a fully
observed Phase-I set $\calIone$ of size $N_1$ and a deployment set
$\calItwo$.  For LPB construction, define the acquisition envelope
\begin{equation}
    Q_i:=\qprior(X_i)=\fhat_{\tauprior}(X_i),
\end{equation}
whereas for metric estimation we set $Q_i=\maxl$.  Phase I observes
\begin{equation}
    L_i:=\min(T_i,Q_i)
\end{equation}
interactions for each $i\in\calIone$.  Its realized cost is
$B_1=\sum_{i\in\calIone}L_i$, leaving the nominal Phase-II budget
\begin{equation}
    \bar B_2=\frac{N\bar B-B_1}{N-N_1},
    \qquad \bar B:=B/N.
    \label{eq:generalized_phase2_budget}
\end{equation}

At every active prefix, we compute the causal hazard score
\begin{equation}
    S_i(t)
    :=\widehat{\mathbb P}
    (T_i=t\mid T_i\geq t,H_i(t),\Dtrain).
    \label{eq:causal_hazard_score}
\end{equation}
This score uses $X_{it}$, represented here by the full history $H_i(t)$, not
only $X_{i0}$.  Importantly, $S_i(t)$ is available before deciding whether to
acquire interaction $t$.  Future histories $H_i(t+1),H_i(t+2),\ldots$ are
never used in that decision.

The direct-bin variants used by Definitive and Generalized DAPRO use $K$
score bins at each time; the default is $K=2$.  On the policy-fitting rows
active at time $t$, compute empirical score quantiles at
$1/K,\ldots,(K-1)/K$.  These cutpoints define a bin
$k_i(t)\in\{1,\ldots,K\}$ for every fitting and deployment score.  Their
policy has the compact form
\begin{equation}
    P_i(t)=p_{t,k_i(t)}.
    \label{eq:score_bin_policy}
\end{equation}
Empty fitting bins copy the nearest observed bin, and the resulting table is
made nondecreasing in the score bin.  Thus, a larger estimated event hazard
never receives a smaller continuation probability at the same time.

\subsection{Target-$A$ DAPRO}

Target-$A$ DAPRO aligns the hard Phase-I objective with a fixed binary target.
For LPB construction, choose an anchor quantile level $\tau_\star<\tauprior$.
The implementation supports two anchors.  The \emph{raw-$\alpha$} anchor
selects the last candidate in the initial quantile grid whose nominal
miscoverage is strictly below the target $\tautarget$.  The
\emph{Phase-I-unweighted} anchor computes the empirical unweighted
miscoverage of every candidate on the fully observed Phase-I rows and applies
the same strict initial-prefix selector.  In either case the selected anchor
is frozen before Phase II; the main Definitive and Generalized variants use
the raw-$\alpha$ anchor with $\tautarget=0.10$.  Define
\begin{equation}
    A_i^{\mathrm{LPB}}
    :=\ind{T_i<\fhat_{\tau_\star}(X_i)}.
    \label{eq:lpb_target_a}
\end{equation}
The strict inequality matches the LPB miscoverage convention.  For metric
estimation, the anchor is replaced by the fixed horizon,
\begin{equation}
    A_i^{\mathrm{metric}}:=\ind{T_i\leq\maxl}.
    \label{eq:metric_target_a}
\end{equation}
The corresponding hard endpoint coefficients are
\begin{equation}
    a_i^{\mathrm{hard}}(t)
    :=A_i\ind{t=L_i},
    \qquad
    d_i^{\mathrm{hard}}(t):=\ind{t\leq L_i}.
    \label{eq:hard_dapro_coefficients}
\end{equation}
Substituting~\eqref{eq:hard_dapro_coefficients} into
\eqref{eq:generalized_dapro_problem} gives
\begin{equation}
    \frac{1}{N_1}\sum_{i\in\calIone}
    A_i\left(\frac{1}{\rho_i(L_i)}-1\right),
\end{equation}
which is the empirical analogue of the exact variance in
Theorem~\ref{thm:metric_ht_variance}.

In the original Target-$A$ implementation, the empirical problem first
produces per-prefix oracle probabilities on Phase I and then projects them
onto a deployable score-to-probability map, such as a time-specific monotone
regression.  Thus, ``Target-$A$'' identifies the hard target-weighted
objective; it does not, by itself, prescribe the projection architecture.
The direct-bin architecture below removes this second regression stage and
optimizes the deployable lookup table itself.

The target anchor must be fixed using Phase I before Phase-II acquisition.
Selecting it after observing the weighted Phase-II outcomes would feed the
acquisition realization back into the target and invalidate the simple
design-based argument.

\subsection{Definitive DAPRO: regularization and direct deployment}

Hard Target-$A$ objectives can be sparse.  If only a few Phase-I rows satisfy
$A_i=1$, the empirical optimizer may saturate those rows and be indifferent
to all remaining prefixes, leaving useful budget unspent and producing a
brittle policy.  Definitive DAPRO stabilizes the target weight with
\begin{equation}
    \widetilde A_i
    :=\frac{A_i+\delta_{\mathrm{glob}}}
    {1+\delta_{\mathrm{glob}}},
    \qquad \delta_{\mathrm{glob}}=0.001,
    \label{eq:definitive_regularization}
\end{equation}
and replaces $A_i$ by $\widetilde A_i$ in
\eqref{eq:hard_dapro_coefficients}.  The small global component assigns a
well-defined value to improving reach outside the currently observed target
events while preserving the target emphasis.

Definitive DAPRO also fixes the deployable class to the direct two-bin policy
in~\eqref{eq:score_bin_policy}.  In contrast to a per-time Platt or isotonic
regression fitted to an unconstrained oracle probability matrix, the bin
table is itself optimized and is deployed by deterministic lookup.  This
removes an unnecessary score-to-probability regression layer and makes the
policy-transfer error easier to diagnose.

Accordingly, Definitive DAPRO is not merely Target-$A$ DAPRO with another
definition of $A_i$.  It is a particular stabilized instantiation: a raw
$\tautarget$ anchor for LPBs or a fixed $\maxl$ anchor for metrics, global
regularization, a small direct score-bin class, a positive deployment floor,
and an explicit budget-transfer mechanism.

There are two budget-control realizations of this same hard objective.
\emph{Definitive DAPRO} uses the projection margin described below, whereas
\emph{Definitive DAPRO+CRC} uses an independent control fold.  The public
\texttt{DAPRO} implementation denotes the latter, finite-sample-controlled
version.  Neither controller changes $\widetilde A_i$ or the optimized
score-bin policy class.

\subsection{Generalized DAPRO: soft event mass at every prefix}

Generalized DAPRO retains the same Phase-I split, causal score, bin table,
sequential optimizer, and deployment mechanism as Definitive DAPRO.  It
changes only the estimator of the target coefficients.  Rather than placing a
noisy one-hot mass at the realized endpoint, it uses the conditional event
mass available at every observed prefix:
\begin{equation}
    h_i(t):=S_i(t)
    =\widehat{\mathbb P}(T_i=t\mid T_i\geq t,H_i(t),\Dtrain).
\end{equation}
For LPB construction, the soft coefficient is
\begin{equation}
    a_i^{\mathrm{soft,LPB}}(t)
    :=\ind{t\leq L_i}
    \frac{
      h_i(t)\ind{t<\fhat_{\tau_\star}(X_i)}
      +\delta_{\mathrm{glob}}h_i(t)
    }{1+\delta_{\mathrm{glob}}}.
    \label{eq:soft_lpb_mass}
\end{equation}
For unsafe-event-rate estimation,
\begin{equation}
    a_i^{\mathrm{soft,metric}}(t)
    :=h_i(t)\ind{t\leq L_i}\ind{t\leq\maxl},
    \qquad
    d_i^{\mathrm{soft}}(t):=\ind{t\leq L_i}.
    \label{eq:soft_metric_mass}
\end{equation}
Because the metric target and acquisition envelope are both $\maxl$, every
active prefix lies inside the named target; consequently the global
regularization in~\eqref{eq:soft_lpb_mass} algebraically cancels for this
particular metric target.

Equation~\eqref{eq:soft_metric_mass} is a Rao--Blackwellization of the hard
endpoint mass: for a row still active at $t$,
\begin{equation}
    \E[\ind{T_i=t}\mid H_i(t),T_i\geq t]=h_i(t).
\end{equation}
It spreads objective mass over plausible event times, substantially reducing
the sampling noise of the Phase-I optimization while preserving causal
deployment.  If one instead sets
$a_i(t)=A_i\ind{t=L_i}$, Generalized DAPRO reduces exactly to the hard
Target-$A$ objective.  In this precise sense, the hard Target-$A$ and
Definitive objectives are special cases of the generalized event-mass
formulation.

\subsection{Solving the score-monotone sequential problem}

Write $y_i(t):=\log P_i(t)\leq0$.  Then
$\rho_i(t)=\exp(\sum_{s\leq t}y_i(s))$.  For fixed score-bin assignments, the
objective and budget in~\eqref{eq:generalized_dapro_problem} become
\begin{equation}
\begin{split}
    \mathcal{L}(y;\lambda)
    ={}&\frac{1}{n}\sum_{i,t}a_i(t)
      \exp\!\left(-\sum_{s\leq t}y_i(s)\right)\\
    &+\lambda\left\{
      \frac{1}{n}\sum_{i,t}d_i(t)
      \exp\!\left(\sum_{s\leq t}y_i(s)\right)
      -\bar B_{\mathrm{fit}}
    \right\}.
    \label{eq:dapro_log_lagrangian}
\end{split}
\end{equation}
At each time, equal score-bin identifiers are tied and continuation
probabilities are constrained to be nondecreasing in the score.  Coordinate
descent updates a complete time column.  The one-dimensional monotonic update
is solved by generalized pool-adjacent-violators (PAV); adjacent score blocks
are merged whenever their unconstrained updates violate monotonicity.  An
outer geometric bisection over $\lambda$ finds a budget-feasible solution.
The implementation uses at most $60$ outer iterations and $30$ inner
coordinate iterations, with log probabilities clipped at $-700$ for numerical
stability.

After solving, each time/bin cell receives the mean fitted continuation
probability of its active Phase-I members.  Unobserved bins copy the nearest
observed bin.  This table is used directly at deployment; no Phase-II event
time is used to define a score cutpoint or continuation probability.

\subsection{Positivity and conversion between cumulative and conditional probabilities}

The raw optimizer may assign very small terminal propensities.  To preserve
finite HT weights, the deployed cumulative reach is mixed with an
always-continue exploration policy.  If $\rho_i^{\mathrm{raw}}(t)$ is the raw
cumulative reach and $\eps_{\pi}=0.005$, define
\begin{equation}
    \rho_i^{\mathrm{floor}}(t)
    :=\eps_{\pi}+(1-\eps_{\pi})
      \rho_i^{\mathrm{raw}}(t).
    \label{eq:terminal_floor_mixture}
\end{equation}
The conditional continuation probabilities deployed by the simulator are
then recovered exactly as
\begin{equation}
    P_i(1)=\rho_i(1),
    \qquad
    P_i(t)=\frac{\rho_i(t)}{\rho_i(t-1)},\quad t\geq2.
    \label{eq:cumulative_to_conditional}
\end{equation}
Unlike clipping only the final product, this construction produces a coherent
sequential policy and guarantees every terminal propensity is at least
$\eps_{\pi}$.

\subsection{Projection-based budget control}

The no-CRC variants use a projection assumption.  Let $\eta\geq0$ be a
budget-transfer margin; the default is $\eta=1$ interaction per Phase-II row.
The Phase-I policy shape is fitted to
\begin{equation}
    \bar B_{\mathrm{fit}}:=\bar B_2-\eta.
\end{equation}
After the bin policy and positivity mixture are constructed, a single shared
logit intercept is applied to the cumulative probabilities.  Bisection
selects the largest intercept whose Phase-I expected cost is at most
$\bar B_{\mathrm{fit}}$.  The same intercept and bin table are then applied
unchanged to Phase II.

\begin{proposition}[Budget validity under projection accuracy]
\label{prop:dapro_projection_budget}
Suppose the expected Phase-II cost of the frozen projected policy, conditional
on the policy-fitting data, is at most its corrected Phase-I expected cost
plus $\eta$.  Then the expected total cost is at most $B=N\bar B$.
\end{proposition}

\begin{proof}
The corrected Phase-I expected policy cost is at most $\bar B_2-\eta$.
The assumed transfer bound therefore makes the Phase-II expected cost at most
$\bar B_2$.  Adding the realized fully observed Phase-I cost $B_1$ and using
\eqref{eq:generalized_phase2_budget} gives
$B_1+(N-N_1)\bar B_2=B$.
\end{proof}

The common intercept guarantees the fitted-sample constraint; it does not by
itself establish the transfer premise of
Proposition~\ref{prop:dapro_projection_budget}.  This is why an independent
budget-control fold remains valuable even after the numerical optimization
has met its target exactly.

\subsection{Independent CRC budget control}
\label{sec:generalized_dapro_crc}

CRC changes the budget controller, not the statistical target.  Applied to
the regularized hard endpoint masses, it gives Definitive DAPRO+CRC; applied
to the soft prefix masses, it gives Generalized DAPRO+CRC.  In either case,
the CRC variant divides the $N_1$ fully observed rows into a policy-fitting
set $\calIone^{\mathrm{fit}}$ of size $N_1-n_{\mathrm{crc}}$ and an
independent budget-control set $\calIone^{\mathrm{crc}}$ of size
$n_{\mathrm{crc}}$.  In our experiments,
$n_{\mathrm{crc}}=N_1/2$.  Only the fitting set determines the target anchor,
score cutpoints, bin policy, and base cumulative reach.  The event times in
the CRC set are revealed only after this candidate family is frozen.

Let $\bar\rho_i(t)$ denote the base cumulative policy.  Before CRC selection,
we optionally cap each row's label-free horizon cost by
\begin{equation}
    C_{\max}:=\min(\maxl,2\bar B).
\end{equation}
If $\sum_{t=1}^{Q_i}\bar\rho_i(t)>C_{\max}$, mix that row toward the floor:
\begin{equation}
    \bar\rho_i^{\mathrm{cap}}(t)
    =\eps_\pi+\kappa_i
      (\bar\rho_i(t)-\eps_\pi),
\end{equation}
where $\kappa_i\in[0,1]$ is the largest value whose horizon cost does not
exceed $C_{\max}$.  This cap is label-free because it uses $Q_i$, not $T_i$.
It is not the target average budget.  Its purpose is to tighten the bounded
loss envelope required by CRC while still permitting the policy to spend more
than the average budget on a high-risk row.

Construct a fixed nested family with $K_{\mathrm{crc}}=401$ candidates,
\begin{equation}
    \rho_i^{(k)}(t)
    :=\eps_\pi+\gamma_k
    (\bar\rho_i^{\mathrm{cap}}(t)-\eps_\pi),
    \qquad
    1=\gamma_1>\cdots>\gamma_{K_{\mathrm{crc}}}=0.
    \label{eq:dapro_crc_family}
\end{equation}
For a fully observed CRC row, let
\begin{equation}
    b_i:=L_i,
    \qquad
    c_i(k):=\sum_{t=1}^{L_i}\rho_i^{(k)}(t).
\end{equation}
Here, $b_i$ is its realized pilot cost, while $c_i(k)$ is the exact expected
cost that candidate $k$ would incur on the same complete trajectory.

Let $n=n_{\mathrm{crc}}$, $m=|\calItwo|$, and let
\begin{equation}
    B_{\mathrm{rem}}
    :=B-\sum_{i\in\calIone^{\mathrm{fit}}}L_i
\end{equation}
be the budget remaining after policy fitting.  Define
\begin{equation}
    \ell_i(k):=c_i(k)+\frac{n}{m}b_i,
    \qquad
    L_{\max}:=C_{\max}+\frac{n}{m}\maxl.
\end{equation}
The CRC rule selects the first, and hence most aggressive, candidate in the
nested family satisfying
\begin{equation}
    \frac{
      \sum_{i\in\calIone^{\mathrm{crc}}}\ell_i(k)+L_{\max}
    }{n+1}
    \leq \frac{B_{\mathrm{rem}}}{m}.
    \label{eq:dapro_crc_selector}
\end{equation}

\begin{theorem}[Marginal expected-budget validity of DAPRO+CRC]
\label{thm:dapro_crc_budget}
Assume the CRC and deployment rows are exchangeable conditional on the
policy-fitting data; the candidate family in
\eqref{eq:dapro_crc_family} is fixed before the CRC event times are inspected;
the candidates are nested; and $0\leq c_i(k)\leq C_{\max}$,
$0\leq b_i\leq\maxl$.  If the selector in
\eqref{eq:dapro_crc_selector} is feasible, then the selected policy satisfies
\begin{equation}
    \E\!\left[
      \sum_{i\in\calI}\Ttilde_i
      \middle|\Dtrain,\Dcalone^{\mathrm{fit}}
    \right]\leq B,
    \label{eq:dapro_crc_expected_budget}
\end{equation}
where the expectation is marginal over the independent CRC sample,
deployment trajectories, and acquisition draws.
\end{theorem}

\begin{proof}
The loss $\ell_i(k)$ is nested in $k$ and bounded by $L_{\max}$.  The standard
conformal risk control leave-one-out argument
\cite{angelopoulos2024conformal} applied to
\eqref{eq:dapro_crc_selector} yields
\begin{equation}
    m\E[c_{\mathrm{new}}(\widehat k)]
    +n\E[b_{\mathrm{new}}]
    \leq B_{\mathrm{rem}}.
\end{equation}
The first term is the expected deployment cost and the second is the expected
cost of the fully observed CRC fold.  Adding the fixed policy-fitting cost
proves~\eqref{eq:dapro_crc_expected_budget}.
\end{proof}

The guarantee in Theorem~\ref{thm:dapro_crc_budget} is marginal.  It does not
assert that every realized calibration split has conditional expected cost at
most $B$, nor that every realized set of Bernoulli acquisitions consumes at
most $B$.  Those stronger statements are neither required by our expected-
budget formulation nor implied by CRC.

The choice $C_{\max}=2\bar B$ is practical rather than fundamental.  Taking
$C_{\max}=\bar B$ prevents targeted overspending on high-risk rows; taking
$C_{\max}=\maxl$ is valid but makes the CRC correction
$L_{\max}/(n+1)$ unnecessarily large.  The intermediate value $2\bar B$
retains useful allocation flexibility while providing a substantially tighter
finite-sample envelope.

\subsection{Metric-optimal PMF: the pre-run instantiation}
\label{sec:metric_optimal_pmf}

Generalized DAPRO can use $H_i(t)$ because it decides whether to continue only
after prefix $t$ has become available.  A policy that commits all
continuation probabilities before running any conversation cannot use future
$H_i(t)$.  Metric-optimal PMF is the corresponding pre-run instantiation of
\eqref{eq:generalized_dapro_problem}.

From the predictive distribution available at the initial prompt, define
\begin{equation}
    \widehat a_i(t)
    :=\widehat{\mathbb P}(T_i=t\mid X_{i0}),
    \qquad
    \widehat d_i(t)
    :=\widehat{\mathbb P}(T_i\geq t\mid X_{i0}).
    \label{eq:pmf_coefficients}
\end{equation}
Only the current-time-zero row of the model's conditional PMF is used.  Using
a later row would leak a future history into a policy that is supposed to be
fixed before interaction.  Any probability mass beyond $\maxl$ is treated as
survival through the metric horizon.  If the model instead returns discrete
hazards $h_i(t)$, compute
\begin{equation}
    \widehat d_i(1)=1,
    \qquad
    \widehat d_i(t)=\prod_{s<t}(1-h_i(s)),
    \qquad
    \widehat a_i(t)=h_i(t)\widehat d_i(t).
\end{equation}

Metric-optimal PMF solves
\begin{equation}
\label{eq:metric_pmf_problem}
\begin{aligned}
    \min_{\rho}\quad
    &\sum_{i=1}^{N}\sum_{t=1}^{\maxl}
      \frac{\widehat a_i(t)}{\rho_i(t)}\\
    \text{s.t.}\quad
    &\frac{1}{N}\sum_{i=1}^{N}\sum_{t=1}^{\maxl}
      \widehat d_i(t)\rho_i(t)\leq\bar B,\\
    &1\geq\rho_i(1)\geq\cdots\geq\rho_i(\maxl)
      \geq\eps_\pi,\quad i=1,\ldots,N,
\end{aligned}
\end{equation}
with $\eps_\pi=1/\maxl$ by default.  The monotonicity constraints are required
because $\rho_i(t)$ is a cumulative reach probability.

The solution has a simple exact construction.  Ignoring temporal monotonicity
for a moment, the KKT conditions yield
\begin{equation}
    \rho_i(t)
    =\operatorname{clip}
      \left(s\sqrt{\frac{\widehat a_i(t)}{\widehat d_i(t)}},
      \eps_\pi,1\right)
    \label{eq:pmf_unpooled_solution}
\end{equation}
for a common scale $s\geq0$.  To impose nonincreasing reach, apply antitonic
PAV separately to every row.  Begin with singleton time blocks.  For a block
$\mathcal{B}$, set
\begin{equation}
    r_{i,\mathcal{B}}
    :=\sqrt{
      \frac{\sum_{t\in\mathcal{B}}\widehat a_i(t)}
           {\sum_{t\in\mathcal{B}}\widehat d_i(t)}
    }.
    \label{eq:pmf_pav_block}
\end{equation}
Whenever a left block has a smaller value than its right neighbor, merge the
two blocks and recompute~\eqref{eq:pmf_pav_block}.  Expand the final block
values to a nonincreasing base path $r_i(t)$, and set
\begin{equation}
    \rho_i(t;s)=\operatorname{clip}(sr_i(t),\eps_\pi,1).
    \label{eq:pmf_scaled_path}
\end{equation}
A one-dimensional bisection chooses the largest common $s$ satisfying the
model-predicted budget in~\eqref{eq:metric_pmf_problem}.  Finally, convert
$\rho_i(t)$ to conditional continuation probabilities using
\eqref{eq:cumulative_to_conditional}.  The PAV construction is
$\mathcal{O}(N\maxl)$, and the bisection requires only repeated linear-time
budget evaluations.

Metric-optimal PMF does not require a calibration split: the predictive model
is trained separately on $\Dtrain$, and no event time from the evaluation
benchmark is used to choose the schedule.  Conditional on the initial PMFs,
the numerical bisection enforces
\begin{equation}
    \frac{1}{N}\sum_{i,t}\widehat d_i(t)\rho_i(t)\leq\bar B.
\end{equation}
This is a model-based budget statement, not a distribution-free one.  Its
true expected cost is
\begin{equation}
    \frac{1}{N}\sum_{i,t}d_i(t)\rho_i(t),
    \qquad
    d_i(t):=\mathbb P(T_i\geq t\mid X_{i0}),
\end{equation}
and can differ from the target when $\widehat d_i(t)$ is misspecified.  HT
validity is unaffected because it uses the realized, known propensities; only
the budget claim depends on model calibration.

\subsection{Implementation summary}

Algorithm~\ref{alg:generalized_dapro_implementation} gives an implementation-
level summary of Generalized DAPRO.  Replacing the soft masses in line~7 by
the hard endpoint masses in~\eqref{eq:hard_dapro_coefficients} gives Target-
$A$ DAPRO; additionally applying~\eqref{eq:definitive_regularization} gives
Definitive DAPRO.  Replacing the LPB target mask by
$\ind{t\leq\maxl}$ gives the metric-estimation variant.

\begin{algorithm}[htbp]
\caption{Generalized DAPRO for LPB construction or metric estimation}
\label{alg:generalized_dapro_implementation}
\begin{algorithmic}[1]
\REQUIRE Prompts $\{X_i\}_{i=1}^N$, predictive quantiles, conditional PMF
model, total budget $B=N\bar B$, $N_1$, target type, $K=2$ bins,
$\delta_{\mathrm{glob}}=0.001$, and $\eps_\pi=0.005$.
\STATE Set $Q_i\gets\qprior(X_i)$ for LPBs and $Q_i\gets\maxl$ for metrics.
\STATE Randomly select $\calIone$ with $|\calIone|=N_1$; observe each row to
$L_i=\min(T_i,Q_i)$.
\IF{CRC is enabled}
    \STATE Split $\calIone$ into policy-fit and CRC folds, with
    $n_{\mathrm{crc}}=N_1/2$.
\ELSE
    \STATE Use all of $\calIone$ for policy fitting.
\ENDIF
\FOR{every active policy-fit prefix $(i,t)$}
    \STATE Compute $S_i(t)=\widehat{\mathbb P}(T_i=t\mid T_i\geq t,H_i(t))$.
    \STATE Set $a_i(t)$ by~\eqref{eq:soft_lpb_mass} for LPBs or
    \eqref{eq:soft_metric_mass} for metrics, and set $d_i(t)=\ind{t\leq L_i}$.
\ENDFOR
\STATE At every time, estimate $K-1$ empirical score-quantile cutpoints on
active policy-fit rows and assign score bins.
\STATE Solve~\eqref{eq:generalized_dapro_problem} in log conditional
probabilities using coordinate descent, generalized PAV, and an outer dual
bisection.
\STATE Convert the fitted solution to a monotone time/bin lookup table; fill
empty bins from the nearest observed bin.
\IF{CRC is disabled}
    \STATE Reserve projection margin $\eta=1$, apply the cumulative-probability
    floor, and fit one Phase-I budget intercept; deploy the frozen table and
    intercept on $\calItwo$.
\ELSE
    \STATE Apply the row-local horizon cap $C_{\max}=\min(\maxl,2\bar B)$.
    \STATE Form the $401$ nested cumulative policies in
    \eqref{eq:dapro_crc_family}; evaluate $b_i$ and $c_i(k)$ only on the CRC
    fold; select $\widehat k$ by~\eqref{eq:dapro_crc_selector}.
    \STATE Deploy the selected frozen policy on $\calItwo$.
\ENDIF
\FOR{each deployment row and each active time $t$}
    \STATE Observe $H_i(t)$, compute its score and bin, look up $P_i(t)$, draw
    $Z_i(t)\sim\operatorname{Ber}(P_i(t))$, and stop if $Z_i(t)=0$, an event
    occurs, or $Q_i$ is reached.
\ENDFOR
\STATE Record the terminal propensity $\pi_i=\prod_{t=1}^{\min(T_i,Q_i)}P_i(t)$
for every deployment row; use $\pi_i=1$ on fully observed rows.
\RETURN Censoring times, exact pathwise propensities, and all data required by
the LPB calibrator or HT metric estimator.
\end{algorithmic}
\end{algorithm}

For clarity, Table~\ref{tab:dapro_variants} summarizes the deployable variants
without the historical legacy objective.

\begin{table}[htbp]
\centering
\caption{Unified view of the deployable allocation variants.  ``Hard'' means
one realized endpoint mass; ``soft'' means conditional event mass at every
observed prefix.}
\label{tab:dapro_variants}
\small
\begin{tabularx}{\linewidth}{lXXXX}
\toprule
Method & Target coefficients & Policy information & Policy class & Budget statement\\
\midrule
Target-$A$ DAPRO
& Hard $A_i\ind{t=L_i}$
& $H_i(t)$ through $S_i(t)$
& Score-monotone projected policy
& Projection-transfer assumption\\
Definitive DAPRO
& Regularized hard endpoint
& $H_i(t)$ through hazard score
& Direct $K=2$ score-bin table
& Margin-based projection assumption\\
Definitive DAPRO+CRC
& Same regularized hard endpoint
& Current causal history $H_i(t)$
& Same direct table plus nested contraction
& Distribution-free marginal expected-total-budget guarantee\\
Generalized DAPRO
& Soft prefix hazards
& Current causal history $H_i(t)$
& Direct $K=2$ score-bin table
& Margin-based projection assumption\\
Generalized DAPRO+CRC
& Same soft prefix hazards
& Current causal history $H_i(t)$
& Same table plus nested contraction
& Distribution-free marginal expected-total-budget guarantee\\
Metric-optimal PMF
& Initial-PMF event mass
& Initial prompt $X_{i0}$ only
& Row-specific precommitted cumulative path
& Exact model-predicted budget; true budget requires PMF calibration\\
\bottomrule
\end{tabularx}
\end{table}

This unified perspective explains why the same framework is useful for both
predictive-bound construction and metric estimation.  The sequential
optimization machinery, causal deployment, positivity correction, and budget
controllers remain unchanged.  What changes is the statistical quantity
whose squared influence is placed in $a_i(t)$.  For an LPB, this mass targets
the miscoverage event below a fixed candidate bound.  For the unsafe event
rate, it targets events by $\maxl$.  For another smooth population metric, the
same machinery can be reused after replacing the event mass by an estimate of
the corresponding squared influence contribution.
